\documentclass[10pt,a4paper]{amsart}
\usepackage[margin=19mm]{geometry}
\usepackage{amsmath,amssymb,mathtools}
\usepackage[scr=rsfs,cal=euler]{mathalfa}
\usepackage{xfrac}
\usepackage[unicode,hidelinks,bookmarks=false]{hyperref}
\newcommand{\ie}{i.e.\ }
\newcommand{\cf}{cf.\ }
\newcommand{\resp}{respectively\ }
\newcommand{\Subsection}{\subsection}
\providecommand{\nobreakdash}{\nobreak\hbox{-}}
\setlength{\emergencystretch}{2em}
\multlinegap0pt
\usepackage{bm}
\usepackage{bbm}
\usepackage{dsfont}
\usepackage{xspace}
\usepackage{cancel}
\usepackage{mathtools}
\usepackage{amsthm}
\makeatletter
\@ifundefined{lem}{\newtheorem{lem}{Lemma}}{}
\@ifundefined{prop}{\newtheorem{prop}{Proposition}}{}
\makeatother
\title[Elliptic fibrations on toric $K3$ hypersurfaces and mirror s]{Elliptic fibrations on toric $K3$ hypersurfaces and mirror symmetry derived from Fano polytopes}
\author{Tomonao Matsumura, Atsuhira Nagano}
\date{Source: arXiv:2208.01465v2 (CC BY 4.0)}
\begin{document}
\maketitle

\noindent\emph{Source and licence.} Elliptic fibrations on toric $K3$ hypersurfaces and mirror symmetry derived from Fano polytopes. Version arXiv:2208.01465v2 (CC BY 4.0), distributed under the Creative Commons Attribution 4.0 International licence, as recorded in the arXiv licence metadata for this exact version and shown on the abstract page https://arxiv.org/abs/2208.01465v2. Full rights notice: licenses/arxiv-2208.01465-CC-BY-4.0.txt.\par\smallskip

\noindent\textit{Editorial scope.} Counted unit: the Lem whose source label is \texttt{LemFinAbel} in the article (source0.tex lines 1836--1861, source label \texttt{LemFinAbel}), delivered together with the article's own complete original proof of it (source0.tex lines 1863--1880, one \texttt{proof} environment of 18 lines). No mathematics is written, repaired or re-derived: statement and proof are the article's own text line for line. Required context retained uncounted: PropEvident (lines 1630--1639). Every definition, display and label of those lines is retained unchanged. Omitted from this package and not redistributed: 1--1629, 1640--1835, 1862--1862, 1881--2941. Nothing inside a retained excerpt is omitted or altered: every retained body is delivered exactly as the article's lines are written, so no omission receipt of any kind accompanies this package. Citations: the retained excerpts carry no citation command at all, so no bibliography entry of the article is transcribed and no reference is redistributed. Reference adaptation: none was needed and none was made; every reference inside the delivered unit resolves inside it, so no unresolved reference or citation is produced. Printed slips kept literal: no printed word, symbol, hypothesis or number of the counted statement or of its proof was altered. Layout and class: the article's own class is \texttt{article} with packages including mathrsfs, bm, graphics, amsfonts, amsthm, amssymb, amsmath, amscd, booktabs, url, longtable, tikz; the wrapper is the lane's pinned \texttt{amsart} editorial wrapper at 10pt on A4, which restates the theorem-like environments the retained text needs and adds the article's own macro definitions that the retained text needs (none), with extra wrapper packages (bm, bbm, dsfont, xspace, cancel, mathtools). The delivered numbering is the unit's own. Assets: no retained span contains a figure environment, a graphics-inclusion command, a PDF-page inclusion, a table or a TikZ picture; no image file is copied into this package, the package declares no asset dependency and no image file is redistributed with it. Licence: version arXiv:2208.01465v2 of this article is distributed under the Creative Commons Attribution 4.0 International licence, as recorded in the arXiv licence metadata for this exact version and shown on the abstract page https://arxiv.org/abs/2208.01465v2. A full rights notice is delivered at licenses/arxiv-2208.01465-CC-BY-4.0.txt. Characters: every retained excerpt is pure ASCII, so the delivered engine, XeLaTeX with its default Latin Modern text font, reports no missing character.
\begin{prop}\label{PropEvident}
For each $k\in \{6,\ldots, 18\},$
let $\ell_k$ be the number of  vertices of the Fano polytope  $P_k$,
 $L_k$ be the lattice given in Table 1
 and $S_k$ be a very general member of the family $\mathscr{F}_k$.
Then, there is a sublattice $E_k$ of ${\rm NS}(S_k)$,
which we call the evident lattice of  $S_k$,
satisfying ${\rm rank} (E_k) =23 - \ell_k$ and $|\det (E_k)| = |\det(L_k)|.$
The generators of $E_k$ are explicitly  given in Table 6.
\end{prop}

\begin{lem}\label{LemFinAbel}
For $k\in\{6,\ldots, 18\},$ let $L_k$ be the lattice of Table 1 and $E_k$ be the evident lattice of Proposition \ref{PropEvident}.
The finite abelian group $\mathscr{A}_{E_k}$ is isomorphic to the group $\mathscr{A}_{L_k}.$ 
Precisely,  setting  
\begin{align*}
& 
\mathcal{G}_6=(\mathbb{Z}/4\mathbb{Z})\oplus (\mathbb{Z}/2\mathbb{Z})^{ \oplus 2} ,\quad 
\mathcal{G}_7=\mathbb{Z}/12 \mathbb{Z}, \quad
\mathcal{G}_8 = \mathbb{Z}/10 \mathbb{Z}, \quad 
\mathcal{G}_9 =\mathbb{Z}/16 \mathbb{Z}, \quad
\mathcal{G}_{10}=\mathbb{Z}/14 \mathbb{Z},
\\
&
\mathcal{G}_{11}=\mathbb{Z}/12 \mathbb{Z}, \quad
\mathcal{G}_{12}=\mathbb{Z}/18 \mathbb{Z}, \quad
\mathcal{G}_{13}=(\mathbb{Z}/14 \mathbb{Z})\oplus (\mathbb{Z}/2 \mathbb{Z}), \quad
\mathcal{G}_{14}=\mathbb{Z}/23 \mathbb{Z},\quad
\mathcal{G}_{15}=\mathbb{Z}/31 \mathbb{Z},\\
&
\mathcal{G}_{16}=(\mathbb{Z}/10\mathbb{Z})\oplus (\mathbb{Z}/2 \mathbb{Z}),\quad
\mathcal{G}_{17}=(\mathbb{Z}/12 \mathbb{Z}) \oplus (\mathbb{Z}/2 \mathbb{Z})^{\oplus 2}, \quad
\mathcal{G}_{18}=\mathbb{Z}/44 \mathbb{Z},
\end{align*}
one has $\mathscr{A}_{E_k} \simeq \mathscr{A}_{L_k} \simeq \mathcal{G}_k$ for each $k$.
Also,  $q_{E_k} \simeq - q_{L_k}$ holds.
\end{lem}

\begin{proof}
For each $k\in\{6,\ldots, 18\},$
let $\{v_1,\ldots,v_{\ell_k-3}\}$ ($\{w_1,\ldots,w_{23-\ell_k}\}$,  resp.)  be the basis of $L_k$  ($E_k$, resp.)
of the intersection matrix of Table 1 
(Table 6, resp.),
where $\ell_k$ is the number of vertices of $P_k$.
Then, we can find an element 
$\alpha \in L_k^\vee / L_k$ 
($\beta\in E_k^\vee/E_k$, resp.) 
in the form
$\alpha = \sum_{i=1}^{\ell_k-3} r_i v_i $
($\beta = \sum_{i=1}^{23-\ell_k} s_i w_i $),
where $r_i$ ($s_i$, resp.) are elements of $\mathbb{Q}/\mathbb{Z}$,
such that $\alpha $ ($\beta$, resp.) generates a cyclic group which is a direct summand of $\mathcal{G}_k$.
One can find the coefficients $r_i$ ($s_i$, resp.) in Table A.2.1 (Table A.2.2, resp.) of Appendix.
The values $q_{L_k}(\alpha)$ and $q_{E_k}(\beta)$ for each $\alpha$ and $\beta$ are listed in Appendix.
Thus, we can check $q_{E_k} \simeq - q_{L_k}$ for every $k\in \{6,\ldots, 18\}$.
\end{proof}

\end{document}
